\begin{table}[H]
\centering
\caption{Distributional Incidence of Liberation Day Tariffs by Income Quintile}
\label{tab:retail_incidence}
\begin{tabular}{lcccc}
\toprule
Income Quintile & Goods Budget & Price Burden & Price Burden & Ratio to \\
& Share & (No Retaliation) & (Retaliation) & Q5 \\
\midrule
Q1 (Lowest) & 41\% & 8.40\% & 5.18\% & 1.41x \\
Q2 & 37\% & 7.58\% & 4.68\% & 1.28x \\
Q3 (Middle) & 35\% & 7.17\% & 4.43\% & 1.21x \\
Q4 & 31\% & 6.35\% & 3.92\% & 1.07x \\
Q5 (Highest) & 29\% & 5.94\% & 3.67\% & 1.00x \\
\midrule
Regressivity ratio (Q1/Q5) & \multicolumn{4}{r}{1.41x} \\
\midrule
\multicolumn{5}{l}{\textit{Illustrative product-level pass-through (retail\_prices\_illustrative.csv):}} \\
\quad Clothing & \multicolumn{4}{r}{30.2\%} \\
\quad Electronics & \multicolumn{4}{r}{29.5\%} \\
\quad Footwear & \multicolumn{4}{r}{29.0\%} \\
\quad Grocery & \multicolumn{4}{r}{30.2\%} \\
\quad Home Appliances & \multicolumn{4}{r}{29.4\%} \\
\bottomrule
\end{tabular}
\begin{tablenotes}
\small
\item Goods budget shares: BLS CEX 2023 Table 1101. Product pass-through computed from retail\_prices\_illustrative.csv.
\end{tablenotes}
\end{table}